% !TeX root = ../../surface-lagrangian-torus-fibration.tex
% 2026-11-21 / v1.2: 共通 anchor による全保持 neck の厳密な moment 範囲．
% 前提：ke:neck-corner-saturated-union と本文中に明示する既存評価．
\section{成分 link の moment cut と固定段階の全域被覆}
\label{ke:component-link-cover-section}

本節では段階 $j$ を先に固定する．$b=b_j$, $n=n_j$, $I=[-\iota_j,\iota_j]$
と書き，系\ref{ke:neck-corner-saturated-union}の部分 fibration を
\[
 f^{\rm old}_{j,t}:U_{j,t}\longrightarrow B^{\rm old}_{j,t}
\]
と表す．以下で $t$ の閾値を小さくすることは許すが，$j$ を $t$ と同時に動かした
一様評価は用いない．Normal cut の新しい深度 $N$ は，$n$ と $\iota_j$ を固定した
\emph{後}に選ぶ．特に $N\gg n$ だけでなく $N\iota_j\gg1$ が必要である．

\subsection{正の residue gauge と actual action periods}

成分 $H_i$ の complement を $W_i=H_i\setminus D_i$ とし，
$D_i=\bigcup_{a=1}^{m}D_a$, $m=d-1$ とする．
正の数 $w_a$ を $\sum_aw_a=m-3$ となるように選ぶ．
$D_m$ を無限遠直線とした affine coordinates の体積形式を
$\Omega_{\rm aff}$，有限直線の defining functions を $\ell_a$ と置く．
固定した component graph gauge $E_t$ 上で
\begin{equation}\label{ke:eq:component-positive-gauge}
 \lambda_t^w
 =d^c\log\frac{(E_t^*\omega_t)^2}{|\Omega_{\rm aff}|^2}
   +\gamma_w,
 \qquad
 \gamma_w=\sum_{a<m}(1-w_a)d^c\log|\ell_a|^2
\end{equation}
と定める．以下では必要に応じて $E_t$ で actual 空間へ押し出して同じ記号を使う．
$\gamma_w$ は $W_i$ 上の固定した閉形式である．
同じ gauge の $t=0$ primitive の meridian period は
$\rho_a:=2\pi w_a$ に対して $-\rho_a$ へ収束する．
本節の全 circle generators は周期 $1$ とする．

\begin{lemma}[Actual meridian period の厳密な規格化]
\label{ke:component-meridian-normalization}
固定段階の component 側 collars と，それらの averaging trajectories を含む
一つの固定 compact の上で\eqref{ke:eq:component-positive-gauge}を用いる．
$\Gamma_a$ は $D_a$ の正の meridian を標識した actual circle とする．
Corner torus の場合には同じ標識の cycle を用いる．
\[
 M_a=\int_{\Gamma_a}\lambda_t^w,\qquad q_a=M_a+\rho_a
\]
と置く．Neck の normal generator の向きと正の meridian の向きの差を
$\epsilon_a\in\{1,-1\}$ とすると，
\begin{equation}\label{ke:eq:component-exact-action-normalization}
 q_a=\epsilon_a\widehat\mu_N
 \quad\text{on a neck collar},\qquad
 (q_a,q_c)=(\nu_a,\nu_c)
 \quad\text{on a component corner collar}.
\end{equation}
ここで $\widehat\mu_N$ は既存の actual canonical normal moment，
$\nu_a,\nu_c$ は小さい triple primitive を全平均して得る moments である．
Corner の source primitive は，非正則な座標補間帯でも必ず
$G^*\lambda_{\log,t}$ として定義する．

これらの等式は加法定数までの等式ではない．また $q_a$ は既存 fibration の全葉上で
一定であり，neck--corner overlap では同じ全 torus 上の同じ function を与える．
独立に先指定した残余 cycle の period が一致する，という主張は含まない．
\end{lemma}

\begin{proof}
まず平均前の形式を同じ正則 graph gauge に置く．有限直線二本が交わる chart で
局所 defining functions を $z_a,z_c$ とすれば，$t=0$ では
\[
 \frac{\Omega_{\log,0}}{\Omega_{\rm aff}}
   =\frac{u_0}{z_az_c}
\]
であり，$u_0$ は corner を跨いで非零の正則 unit である．
固定した punctured collar 上では，$t$ を十分小さくすると
\[
 \frac{E_t^*\Omega_{\log,t}}{\Omega_{\rm aff}}
   =\frac{u_t}{z_az_c},\qquad u_t\longrightarrow u_0
\]
となる．$u_t$ 自体の $z_a=z_c=0$ への延長を仮定する必要はない．
固定 collar 上の非零関数の homotopy により，$u_t$ の二つの meridian winding は
厳密に零である．$d^c\log|z_a|^2$ の正の meridian period は $2\pi$ なので，
\[
 \int_{\Gamma_a^{\rm bare}}(\lambda_t^w-E_t^*\lambda_{\log,t})
 =2\pi(1-w_a)-2\pi=-\rho_a .
\]
同じ計算は $c$ にも成立する．他の defining functions は局所 unit であり，
その meridian winding は零である．

無限遠直線を含む chart では，$\Omega_{\rm aff}$ の pole order が $3$，
logarithmic form の pole order は $1$ である．従って両 canonical primitives の
\emph{差} $\lambda_{{\rm aff},t}-E_t^*\lambda_{\log,t}$ の無限遠 meridian period
は厳密に $4\pi$ である．$\gamma_w$ のその period は
\[
 -2\pi\sum_{a<m}(1-w_a)=-2\pi(2+w_m)
\]
だから，和は $-2\pi w_m$ となる．これにより全直線について同じ式を得る．
この計算が用いるのは固定した divisor の整数 winding であり，
$t\to0$ の近似から未知の加法定数を零と置いているのではない．

Neck の logarithmic volume form は，component の meridian に対して単純対数極を
持つ．従って同じ一変数の計算が成立する．Base-only frame の変更は normal
meridian period を変えず，系\ref{ke:transition-normal-primitive}の Jacobian unit
も winding を持たない．Source の点を含む residual compact でも，
$\mu\ne0$ の branch 上では，例えば
$d\log x'\wedge dp=-d\log y'\wedge dp$ という表示で logarithmic form が延びる．
閉形式 $\lambda_t^w-\lambda_N$ の meridian period は，連結 collar 内の isotopy に
沿って Stokes の定理により一定である．従って一つの source-free chart で確認した
定数が source を含む全 residual collar で同じになる．

次に actual cycles へ移す．Centered averaging map は
\[
 \Phi^*\lambda=\mathcal A\lambda+dQ
\]
を満たし，relative corner patch の Hamiltonian map は
$\Upsilon^*\lambda-\lambda$ を exact にする．従って
$C_t=G\Upsilon\Phi$ による全 marked cycles の period は source の
$\mathcal A(G^*\lambda_{\log,t})$ の period と厳密に同じである．
一方 $\lambda_t^w-\lambda_{\log,t}$ は閉形式なので，その period は
$G$，Moser flow，Hamiltonian flow による meridian isotopy で変わらない．
これが corner 側の等式を与える．Normal averaging について同じ exact transport を
使えば neck 側の等式を得る．非正則な $G$ に対して $G^*d^c=d^cG^*$ を用いていない．
特に $d^c\log((G^*\omega_t)^2/|\Omega|^2)$ を新しい primitive として作り直さない．

最後に，Lagrangian torus 上では primitive の制限は閉形式である．
その period は cycle の同一 torus 内の代表に依存しない．
Neck の全葉では normal moment が一定であり，corner の全葉では全 moments が一定で
ある．系\ref{ke:neck-corner-saturated-union}は共通 normal cycle を厳密に同定し，
第二 cycle は共通 torus を通じて移送するので，全ての表示は compatible である．
\end{proof}

\paragraph{二つの meridians の整数基底．}
Triple chart を $\zeta_i\zeta_a\zeta_c=t$ とし，component $H_i$ 側では
\[
 V_a=(-1,1,0),\qquad V_c=(-1,0,1)
 \quad\text{in }\ker(\mathbb Z^3\xrightarrow{\sum}\mathbb Z)
\]
を取る．これは整数基底であり，それぞれ二つの component boundary lines の正の
meridian である．残る normal subgroup は $V_c-V_a$ である．
別の residual cycle を選んで $V_c$ を $V_c+kV_a$ に変える場合には，
period vector と residue vector の双方を同じ整数行列で変換する．
片方だけを変えて独立な第二 moment の加法定数を零と置くことはしない．

\subsection{深度比に一様な corner moments と閉じた link}

Corner の component 側で二つの独立深度を $a,c$，従属深度を
$L-a-c$ とする．$m_*=\min\{a,c,L-a-c\}$ と置く．
関連 KE 主稿の \texttt{simplex:prop:dominant-product} と
\texttt{simplex:cor:actual-dominant-product} は，$a/c$ に依存しない全固定階の
product-cusp 比較を与える．Canonical volume potentials の差に $d^c$ を施し，
周期 $1$ の meridian の product 長さ $O(a^{-1})$, $O(c^{-1})$ を掛けると，
\begin{equation}\label{ke:eq:component-two-action-estimate}
 \begin{split}
 \nu_a&=\frac{4\pi}{a}
   +O\!\left(\frac{m_*^{-1}+((a+c)/L)^2}{a}\right),\\
 \nu_c&=\frac{4\pi}{c}
   +O\!\left(\frac{m_*^{-1}+((a+c)/L)^2}{c}\right)
 \end{split}
\end{equation}
を得る．ここで $a/c$ の上下界を仮定しない．
正則な plateau での Jacobian unit の exact 項は全平均 period を変えない．
非正則な補間が実際に変化する帯では該当する二つの小深度は固定比の間にあり，
$G-\id$ の指数的な全固定階評価を上の誤差に吸収できる．
正則 plateau の不均衡な領域では normal action の厳密な同変性と
weighted generator 評価を用いる．従って\eqref{ke:eq:component-two-action-estimate}は
相対化された source の moments にも成立する．

この議論は RZ の smooth-divisor expansion を角へ輸入するものではない．
固定 $a,c$ の後に $t\to0$ とすれば，component limit でも相対誤差
$O(1/\min(a,c))$ が得られる．一方，残余基底の固定 compact 上で片方の深度だけを
深くする領域では，既存の fixed-end moment 評価を用いる．

\begin{lemma}[Actual moment cuts による閉 component link]
\label{ke:closed-component-link}
$j$ を十分大きく固定する．各 component の任意の指定 compact を先に与えてよい．
その後 $N$ を十分大きく選び
\[
 \delta=\frac{4\pi}{N}<\frac{\iota_j}{4},\qquad
 N\gg n_j,\qquad N\iota_j\gg1
\]
とする．$|t|$ を十分小さくすれば，各 component graph の中に滑らかな閉じた
三次元 link $S_{i,j,t}$ と，それを境界とする compact domain $K_{i,j,t}$ が存在する．
指定 compact の graph image は $\operatorname{int}K_{i,j,t}$ に含まれ，異なる
components の $K_{i,j,t}$ は互いに素にできる．

Link とその小さい両側 collar は $U_{j,t}$ に含まれる．Neck 部分では
$q_a=\delta$，corner 部分では二つの moment cuts の滑らかな丸めで定義され，
各 link は $f^{\rm old}_{j,t}$ の全 fibers の和である．同じことは collar の平行な
level hypersurfaces にも成立する．Sources の fixed points は link に含まれない．
Residual simple saddles を含む葉は全体として link に含まれてよい．

さらに $\lambda_t^w$ と exact に異なる primitive が，該当する actual normal
actions と corner $T^2$ actions に不変となるように作られているならば，その
Liouville 場は link に外向きに横断する．この最後の主張だけが primitive の
相対不変化を必要とする．
\end{lemma}

\begin{proof}
残余基底の end 深度を $b_{\rm res}$ と書く．各二重線で
$b_{\rm res}<6.5b_j$ までの固定 residual compact を用い，その少し広い compact
も先に固定する．$q_a=\delta$ は
補題\ref{ke:component-meridian-normalization}により
$\widehat\mu_N=\epsilon_a\delta$ そのものであり，$I$ の内部にある．
Fixed-end moment 評価と normal radial 単調性により，$N$ を大きくした後に $t$ を
小さくすれば，この level は normal 深度 $[N/2,2N]$ の間にただ一つ存在する．
この評価は固定 residual compact の全点と全 phase に一様であり，source の基底点も
含む．$\delta\ne0$ なので circle action はこの level 上で自由である．

二つの component boundary lines が交わる corner では，
$a,c\ge6.2b_j$，$a,c\le3N$，$L-a-c\to\infty$ の領域を用いる．
\eqref{ke:eq:component-two-action-estimate}から，各 normal cut は深度
$[N/2,2N]$ に位置し，$a/c$ が大きい場合にもこの範囲は変わらない．
Moments が同時座標になることも確認する．Source の平均形式を $\kappa$ とし，
標準複素構造を $J_0$ とすると，既存の tame 評価から
$\kappa(V,J_0V)>0$ である．深度座標に対して
\[
 \frac{\partial\nu_r}{\partial s_q}
 =-\frac1{4\pi}\kappa(V_r,J_0V_q),\qquad (s_1,s_2)=(a,c).
\]
従って $D\nu$ の対称部分は負定値である．深度の凸領域上で直線を積分すると
\[
 (\nu(s)-\nu(s'))\cdot(s-s')<0\qquad(s\ne s')
\]
となり，$\nu$ は単射で，Jacobian も非退化である．
特に二つの cut の交点は高々一つである．積 moment への相対的な $C^1$ 比較から，
矩形 $[N/2,2N]^2$ の対向する各 face では，$\nu_r-\delta$ はそれぞれ正と負である．
$(4\pi/a,4\pi/c)$ への線型 homotopy はこの境界符号を保つので，degree は $1$ で
あり，$(\delta,\delta)$ の逆像が存在する．同じ符号に固定相対幅の余裕があるため，
moment image は $(\delta,\delta)$ の固定相対幅の近傍を含む．

$0<\varepsilon<\delta/20$ を取り，滑らかな凸関数
$\chi_\varepsilon\ge|\cdot|$ を
$\chi_\varepsilon(r)=|r|$ for $|r|\ge\varepsilon$，
$|\chi_\varepsilon'|\le1$ となるように選ぶ．
\[
 F_{ac}(q_a,q_c)
 =\frac{(q_a-\delta)+(q_c-\delta)
       -\chi_\varepsilon(q_a-q_c)}2
\]
とし，corner の link を $F_{ac}=0$，core 側を $F_{ac}\ge0$ とする．
$\partial_{q_a}F_{ac},\partial_{q_c}F_{ac}\ge0$ で，和は $1$ だから
$dF_{ac}\ne0$ であり，滑らかな丸めになる．
丸めが変化するのは両 moment が $\delta+O(\varepsilon)$ になる領域であり，
そこで二つの深度はいずれも $N$ と比較可能である．

Neck--corner の共通帯 $5b_j<b_{\rm res}<9b_j$ にある link では，一方の
moment が $\delta$，他方は\eqref{ke:eq:component-two-action-estimate}により
$4\pi/b_{\rm res}$ と比較可能である．$N/b_j$ を十分大きく選べば後者は
$\delta+2\varepsilon$ より大きい．従ってこの開帯では
\[
 F_{ac}=q_a-\delta
\]
が\emph{厳密に}成立する．補題\ref{ke:component-meridian-normalization}の
period 同定により，これは neck の同じ actual normal moment cut である．
同じ等式は小さい平行 levels でも成立する．
従って cut を物理的な深度で近似的に貼る必要はない．

Neck end は source torus の $B=7b_j$ で切られている．
$B$ と actual reduced graph の順序保存および $O(b_j^{-2})$ の相対誤差から，
例えば $b_{\rm res}<6.5b_j$ の face はその内側に含まれる．
Corner 側では source の最小深度が $6.2b_j$ より大きければ
\[
 r_{\min}:=\min(\ell_r-6b_j)>0.2b_j>1,\qquad r_{\max}<L
\]
なので $r_{\max}/r_{\min}<L$ である．これは full softmax domain の内点条件
そのものである．$6.4b_j<b_{\rm res}<6.5b_j$ を共通帯に用いれば，neck faces と
corner faces は開集合上で接続し，欠けた annulus は残らない．

こうして得る link は，通常の SNC divisor の深度 $N$ における丸めた tubular
boundary と同じ組合せで貼り合わさる．Normal level の一意性，corner の moment
座標の単射性，および averaging maps の embeddings により自己交差はない．
Normal faces は標準 circle bundle の radial level の graphs であり，corner faces は
正規化した log-depth 座標における小さい $C^1$ 変形と滑らかな丸めである．
既存の $G$，Moser map，Hamiltonian patch は恒等写像からの isotopies を持ち，
共通帯では同じ normal action と同じ cut を保って貼り合わされる．
従ってこの変形は SNC の plumbing の incidence と meridian framings を変えず，
得られた閉 hypersurface は標準 tubular boundary に isotopic である．
固定 $N$ に対する全 charts，元の tori，Moser trajectories は一つの
component compact の graph image に含められる．Link は通常の丸めた boundary に
isotopic なので，その内側の compact domain を $K_{i,j,t}$ と定める．
指定した compact は $N$ をさらに大きくしてこの内側に含められる．
異なる components の固定 compact graph images は，$t$ を小さくすれば互いに素で
ある．

Moment functions は全 fibers 上で一定であり，neck--corner の全交わりは
系\ref{ke:neck-corner-saturated-union}により共通の全 fibers の和である．
従って得られた link と平行 levels は全 fibers の和である．
Source の fixed point では $\widehat\mu_N=0$ なので link 上にない．
Residual foliation の saddle は $q_a=\delta$ の正則性を妨げない．

最後に，相対不変 primitive を $\lambda^{\rm ad}$，その Liouville 場を $Z$ と
する．Exact gauge の変更は periods を変えず，不変性から
\[
 \lambda^{\rm ad}(V_a)=M_a=q_a-\rho_a,\qquad Zq_a=M_a.
\]
したがって normal face では $Z(q_a-\delta)=\delta-\rho_a<0$ であり，corner では
\[
 ZF_{ac}
 =\partial_{q_a}F_{ac}(q_a-\rho_a)
  +\partial_{q_c}F_{ac}(q_c-\rho_c)<0.
\]
ここでは必要なら $N$ を大きくして，丸め部分の全 moments を全 $\rho_a$ より
小さくした．Core 側は $F\ge0$ なので，これは外向き横断性である．
\end{proof}

\subsection{固定段階で残る領域の compact 性}

\begin{lemma}[中央の全特異 strata の固定段階被覆]
\label{ke:fixed-stage-stratum-coverage}
十分大きい各固定 $j$ に対し，中央 fiber の特異 locus の $t$ に依存しない ambient
開近傍 $\mathcal O_j$ と $\tau_j>0$ があり，
\[
 X_t\cap\mathcal O_j\subset U_{j,t}
 \qquad(0<|t|<\tau_j)
\]
となる．従って各 component の固定 compact $K^*_{i,j}\Subset W_i$ を選んで
\begin{equation}\label{ke:eq:fixed-stage-compact-complement}
 X_t\setminus U_{j,t}\subset\bigcup_i E_{i,t}(K^*_{i,j})
\end{equation}
とできる．この結論には source を含む全二重線と全三重点が含まれる．
\end{lemma}

\begin{proof}
まず二重線の residual compact を取る．各 end の深度が $6.5b_j$ 以下の compact
を $P^0$ とし，それを end 深度 $6.7b_j$ 未満の少し広い compact $P^1$ の内部に
含める．Source は $P^0$ の内部に含める．
既存の normal circle と centered chain homotopy をそのまま用い，
補題\ref{ke:fixed-stage-source-neck}をさらに深い固定 normal cutoff $M$ に適用する．
$M$ は $n_j$，$\iota_j^{-1}$，残余 buffer の逆数より十分大きく選ぶ．
この補題の source/image buffer は
\[
 K_{t,M}\subset\operatorname{image}\phi_{N,t}^{(M)},\qquad
 (\phi_{N,t}^{(M)})^{-1}(K_{t,M})\subset U_{t,M}
\]
を与え，forward/backward の actual 長さは $C_{P^1}/M$ 以下である．
Normal circle，平均形式の path，centered kernel は元の $n_j$-tube の制限と同じ
なので，両方が定義された点では ODE の一意性により
\[
 \phi_{N,t}^{(M)}=\phi_{N,t}
\]
である．深い cutoff ごとに別の actual action を選び直してはいない．

ここでは fixed-shell moments の\emph{上界}も必要である．
既存の fixed-end 証明は，深度 $M$ の二つの shell において，全 residual points と
全 phase に一様に
\[
 \mu_N=\frac{4\pi}{M}+O_{P^1}(M^{-1-\nu})+o_{t\to0;M}(1),
 \quad
 \mu_N=-\frac{4\pi}{M}+O_{P^1}(M^{-1-\nu})+o_{t\to0;M}(1)
\]
を与える．従って $M$ を固定した後に $t$ を小さくすれば，両 shell の絶対値は
$C_{P^1}/M$ 以下である．$P^1$ は三重点を避けた固定 compact であり，source の
残余点は成分極限の cusp shell 上で特別な特異点ではないので，同じ評価に含まれる．
Normal radial 単調性により $U_{t,M}$ の全点で
\[
 |\mu_N|\le C_{P^1}/M<\iota_j
\]
となる．Source fiber では二軸上の単調性と固定点で $\mu_N=0$ を使えばよい．

Residual 方向の射影下界と $C_{P^1}/M$ の trajectory 長さによって，$P^0$ 上の
点の逆像は $P^1$ 上に留まる．End の reduced leaves の順序保存と
$b_{\rm res}=B(1+O(b_j^{-2}))$ から，$P^1$ は $B=7b_j$ の end cut の内側にある．
従って $K_{t,M}$ の $P^0$ 部分は\emph{既存の，切り詰め済み} neck に含まれる．
座標の小さい固定 polydisc を取れば，中央二重線の $P^0$ 部分の固定 ambient
近傍の $X_t$ との交わりをこの $K_{t,M}$ に入れられる．

三重点では $m=6b_j$ を用いる．Corner map $C_t=G\Upsilon\Phi$ とその逆の
model 正規化座標での変位は $O(b_j^{-2})$，座標接続の寄与は指数的に小さい．
Clipped depth の議論と同じ $|d\log\ell_r|_h\le C$ の評価により，physical
depths が全て $6.3b_j$ より大きい点の逆像では，$j$ を十分大きく取れば
\[
 \min_r\ell_r(C_t^{-1}x)>6.2b_j
\]
となる．必要な inverse image buffer は $\min\ell_r>m/2$ の既存の開領域内に
ある．従って
\[
 \min_r(\ell_r-m)>0.2b_j>1,\qquad
 \max_r(\ell_r-m)<L
\]
であり，逆像は full softmax domain の\emph{内部}に属する．
この議論により，三つの physical depths が $6.3b_j$ より大きい固定 ambient
近傍と $X_t$ の交わりが actual closed corner に含まれる．

Neck 側の $6.5b_j$ と corner 側の $6.3b_j$ の間には正の幅がある．
End coordinate と triple coordinate の固定 frame 差，および $t$ による小さい
差はこの幅に吸収できる．従って上で作った ambient neighborhoods は中央の全二重線
を三重点まで切れ目なく覆う．粗い条件 $\min\ell_r>2m$ だけを用いると
$7b_j$ と $12b_j$ の間に被覆の gap が残り得るため，ここではそれを用いなかった．
有限個の二重線・三重点について閾値を共通化すると $\mathcal O_j$ を得る．

最後に，中央 fiber の $\mathcal O_j$ の外側は，各 $W_i$ の互いに素な固定 compact
の和に含まれる．族は射影的であり，この部分では fixed component graph が存在する．
従って少し広い固定 compact $K^*_{i,j}\Subset W_i$ を取れば，小さい全 $t$ について
$X_t\setminus\mathcal O_j$ はその graph images に含まれる．
これは，もし反例列があれば射影空間内で中央 fiber の regular point へ収束し，
その点の一意な graph 表示に反することからも分かる．
これで\eqref{ke:eq:fixed-stage-compact-complement}を得る．
\end{proof}

\begin{lemma}[全 fibers による core 外側への切詰め]
\label{ke:whole-fiber-core-trim}
補題\ref{ke:closed-component-link}で，
補題\ref{ke:fixed-stage-stratum-coverage}の各 $K^*_{i,j}$ の graph image を
$\operatorname{int}K_{i,j,t}$ に含めるように $N$ を選ぶ．このとき
\[
 X_t=U_{j,t}\cup\bigcup_i K_{i,j,t}.
\]
また
\[
 U^{\rm ext}_{j,t}
 :=U_{j,t}\setminus\bigcup_i\operatorname{int}K_{i,j,t}
\]
は既存 fibration の全 fibers の和であり，その restriction は全旧 patch 間の
saturation と既存の fiber partition を保つ．各 core との交わりは $S_{i,j,t}$
全体であり，これは双方に用いる全 boundary fibers の和である．

小さい平行 moment levels を用いると，境界の両側に開いた共通 collars を取れる．
従って各 core 側の fibration がこの collar 上で旧 fiber partition と一致するように
構成されているならば，有限個の compact bases を貼り合わせて actual 曲面全体の
Lagrangian fibration を得る．本補題自体は最後の相対 filling の存在を主張せず，
その構成は定理\ref{ke:actual-global-fibration}で与える．
\end{lemma}

\begin{proof}
最初の被覆は\eqref{ke:eq:fixed-stage-compact-complement}と $K^*_{i,j}$ の包含から
従う．次に旧 fibration の一つの fiber $L$ を取る．旧 fibers は regular tori，
一重 pinched tori，または saddle 型の連結な fibers である．
$L$ がある $K_{i,j,t}$ の内側と外側の両方に点を持てば，連結性から
$L\cap\partial K_{i,j,t}\ne\varnothing$ となる．
しかし $\partial K_{i,j,t}=S_{i,j,t}$ は全旧 fibers の和だから，この場合
$L\subset S_{i,j,t}$ でなければならず，内外に点があることと矛盾する．
従って各 $L$ は一つの core 内部，その境界，または全 cores の外側のいずれかに
全体として含まれる．これが切詰めの saturation を証明する．

旧 patch の任意の交わりは，既に系\ref{ke:neck-corner-saturated-union}により
共通の全 fibers の和であった．全 fiber 単位の restriction はこの性質を保つ．
新しい境界で別の half-torus や residual circle segment が現れることはない．
平行 moment levels についても同じ議論を適用できるので，小さい開 common collars
を用意できる．相対 filling がその collar partition を実現すれば，有限 compact
bases の閉じた同値関係による貼合せは Hausdorff であり，局所 fibration の記述も
collar 上で一致する．これで条件付きの最後の主張を得る．
\end{proof}

\begin{lemma}[共通 anchor による保持 neck の moment 範囲]
\label{ke:anchor-trimmed-neck-domain}
系\ref{ke:neck-corner-saturated-union}の一つの旧 neck を
$\mathcal N^{\rm old}_{t,D}$，$D=D_{ik}$ と書き，その共通 moment を
$\widehat\mu$ とする．向きは component $i$ 側で
$q_i=-\widehat\mu$，component $k$ 側で
$q_k=\widehat\mu$ となるように取る．
補題\ref{ke:closed-component-link}の link を共通 anchor
$\delta=4\pi/N_j$ で取り，それが囲む domains を $K_i,K_k$ とする．
既存の選択順序の中で $n_j,N_j$ の下限，最後の $t$ の閾値を強くすれば，
\begin{align*}
 S_i\cap\mathcal N^{\rm old}_{t,D}
   &=\{\widehat\mu=-\delta\}\cap\mathcal N^{\rm old}_{t,D},\\
 S_k\cap\mathcal N^{\rm old}_{t,D}
   &=\{\widehat\mu=\delta\}\cap\mathcal N^{\rm old}_{t,D},\\
 \operatorname{int}K_i\cap\mathcal N^{\rm old}_{t,D}
   &=\{\widehat\mu<-\delta\}\cap\mathcal N^{\rm old}_{t,D},\\
 \operatorname{int}K_k\cap\mathcal N^{\rm old}_{t,D}
   &=\{\widehat\mu>\delta\}\cap\mathcal N^{\rm old}_{t,D}.
\end{align*}
他の component の core はこの neck と交わらない．従って共通 anchor で
whole-fiber trim した neck の moment 範囲は正確に $[-\delta,\delta]$ である．

実際の接着に小さい平行 link を用いた場合も，共通 collar 内で基底の分割を
この anchor へ取り直す．平行 link と anchor の間，及び葉族の補間支持を
当該 component の変更 collar に含めると，core と変更 collar の外に残る旧 neck
には常に $|\widehat\mu|\le\delta$ が成り立つ．
\end{lemma}

\begin{proof}
まず link と旧 neck の交わりを全域で確認する．Fixed residual body と各 end の
$b_{\rm res}<6.5b_j$ の部分では，link は同じ actual normal moment の
$q_i=\delta$ 又は $q_k=\delta$ そのものである．
各 end で葉を source torus の残余深度 $B$ によって標識する．
Neck は $B=7b_j$ の葉で切られており，reduced graph の評価と
式\eqref{ke:eq:neck-residual-first-exit}を全 normal trajectory に積分した評価から
\[
 b_{\rm res}=B\bigl(1+O(b_j^{-2})\bigr)
\]
が成り立つ．従って $B\le6.4b_j$ の部分は上の normal-face 領域内にある．
$6.4b_j\le B\le7b_j$ の部分では，K2 の共通 moment strip により二つの
normal 深度は $4n_j$ より大きい．その全 torus は full softmax domain の内部に
あり，共通 plateau の同じ全 corner torus である．ここでは対応する component
の normal face と corner face の period functions が厳密に一致する．

Component $i$ の境界をこの end で見ると，独立深度の一つは $N_j$ と固定比の間，
残る独立深度は $B(1+O(b_j^{-2}))=O(b_j)$ である．後者の meridian moment を
$q_c$ とすると，式\eqref{ke:eq:component-two-action-estimate}より
$q_c\asymp b_j^{-1}$ である．$N_j/b_j$ を先に大きく取れば，
$q_i=\delta$ の固定した小近傍で
\[
 q_c-q_i>\varepsilon,\qquad F_{ic}=q_i-\delta
\]
が厳密に成り立つ．従って rounding はこの全 overlap では normal face に退化する．
別の face $q_c=\delta$ と rounding の変化領域では，対応する residual 深度は
$N_j$ と固定比の間にあるため，$B\le7b_j$ の旧 neck に入らない．
Component $k$ 側も同様である．これで最初の二つの等式が従う．

この argument の coverage には余裕がある．実際，normal-face 領域の
$b_{\rm res}<6.5b_j$ と corner 側の $B\ge6.4b_j$ は開帯上で重なる．
異なる二重線の outer tubes は互いに素であり，三重点の非 incident charts も
分離されているので，他の normal face がこの neck 内に現れることはない．
さらに $P_j^+\Subset D\setminus Q$ は他の平面 $H_a$，$a\notin\{i,k\}$ と
交わらない．その固定 outer tube の閉包もそれらの平面から離して選べる．
各 $K_a$ は，$N_j$ を固定した後に $t$ を小さく取ると $H_a$ の任意に小さい
近傍に入る固定 compact graph image なので，$K_a$ はこの neck と交わらない．

最後に境界のどちら側が core に属するかを確認する．旧 neck の compact marked
quotient は $[-\iota_j,\iota_j]\times T_{ik}$ であり，その fibers は連結である．
この表示には，end の moment-dependent な cut を順序保存する collar 座標で
自明化したものを用いる．従って
$\widehat\mu< -\delta$ と $\widehat\mu> -\delta$ の各側は連結である．
例えば，これらの各側を compact connected subintervals と $T_{ik}$ の積で
exhaust すれば，proper な connected-fiber map の各逆像は連結である．
この連結性は，逆像が二つの閉集合に分離されたなら，connected fibers と
properness により基底も二つの閉集合に分離されることから直接従う．

両側とも $S_i$ と交わらず，$q_i=\delta$ の直近では
$q_i>\delta$ が core の内側である．従って連結性から
$\widehat\mu<-\delta$ の側全体が $\operatorname{int}K_i$ に入り，反対側は
$K_i$ の外にある．Component $k$ についても同じ議論を使えば残りの二式を得る．

平行 link を用いた場合，anchor における葉族は元の葉族そのものであり，
共通 Liouville collar で anchor へ基底分割を戻せる．この collar と相対補間支持は
$q_a\in[\delta/2,2\delta]$ の一つの固定 component compact 内に含まれる．
それらを変更 collar として扱えば，その外で保持する旧 neck は上で得た
中央 moment 範囲だけである．
\end{proof}

\paragraph{本節の結論と後続の相対 filling．}
本節は，period の加法定数，共通整数基底，二つの moment cuts の丸め，固定段階の
中央特異 locus の被覆，および core 外側への全 fiber 単位の切詰めを，既存 K2 の
maps と quantitative buffers から導いた．Primitive の相対不変化と非捕捉性は
節\ref{ke:core-invariant-primitive-section}，collar 上の葉族の補間は
命題\ref{ke:actual-saddle-moving-neck-interpolation}で証明する．
これらを Liouville flow による相対 filling と合わせるのが
定理\ref{ke:actual-global-fibration}である．
全曲面の ambient metric に関する no-shortcut と normalized distortion は，
本節の固定段階被覆だけからは従わず，
定理\ref{ke:global-ambient-lower}，\ref{ke:global-ambient-upper}で別に評価する．
